\begin{textsample}{NEG Dim 3 – rj – Score 1.00 – rj/rj\_em\_5.txt}  \label{ex:f3_neg_276}
ARQUITETURA DO NOVO ENSINO MÉDIO O Novo Ensino Médio busca atender às exigências sociais e culturais do mundo contemporâneo.
O Currículo Referencial do Rio de Janeiro apresenta o protagonismo no cerne de sua estruturação.
Um currículo que contempla o desenvolvimento integral e o estudante no centro do processo de ensino-aprendizagem, possibilitando aos jovens escolherem itinerários formativos e trilhas de aprofundamento.
O componente Projeto de Vida permite o preparo do educando para a assertividade nas suas escolhas e tomadas de decisão, rumo ao autoconhecimento e criticidade, e também possibilita o fortalecimento da relação dialógica entre professor e estudante, proporcionando um processo avaliativo efetivo, visando o alcance dos objetivos educacionais e a melhoria dos resultados.
Sabe -se que um Currículo é resultado de um conjunto de conhecimentos e saberes selecionados, mas, também, é trajetória, é território e é vivência, que se amplia na relação cotidiana que se estabelece entre professores, estudantes e escola, portanto, enquanto documento de identidade, como bem define Tomaz Tadeu Silva (2009)³, caberá a cada comunidade escolar construir seu Currículo tendo como orientação o Currículo Referencial do Estado do Rio de Janeiro – Ensino Médio ( Regular e Educação de Jovens e Adultos ). ³ SILVA, Tomaz Tadeu.
Documentos de identidade ; uma introdução às teorias do currículo. 3ª ed., Belo Horizonte : Autêntica : 2009.

% matched lemmas: 
\end{textsample}
